% !TeX root = main.tex
% Numerical source: artifacts/runs/R05_v8_261105/<dataset>/validation/diagnosis/
% base_description_graph_graph-activity_graph-role-activity/diagnosis.json, via
% R05_v8_261105/reports/paper_evidence.json. Table rows, contrast cells and prose percentages
% are generated and checked by R05_v8_261105/reports/short_paper_numbers.py (--check).
% The stored 10-pair statistical family is preserved even when only four contrasts are shown.
\section{Experiments}
\label{sec:experiments}
We evaluate HGVRec through three questions: \textbf{RQ1} how it compares with title-only, scene-description, and mean-pooled scene-graph representations; \textbf{RQ2} what changes when ASHA is added to the scene-graph encoder; and \textbf{RQ3} what changes when REAF is added to the encoder that already uses ASHA.

\begin{table}[t]
  \caption{NDCG@10 differences ($\times10^{-3}$), relative change (\%) in parentheses, and Bonferroni-corrected 99.5\% paired-bootstrap intervals.}
  \label{tab:contrasts}
  \centering
  \footnotesize
  \setlength{\tabcolsep}{3pt}
  \begin{tabular}{@{}lcc@{}}
    \toprule
    Contrast & MicroLens-100K & M$^3$L-10M \\
    \midrule
    HGVRec $-$ Title-only & 9.74 (28.90) [8.87, 10.59] & 4.69 (3.38) [3.38, 5.95] \\
    HGVRec $-$ Scene Desc. & 5.51 (14.51) [4.71, 6.30] & 3.75 (2.68) [2.46, 5.01] \\
    +ASHA $-$ Scene Graph & 7.08 (20.14) [6.34, 7.80] & 4.13 (2.99) [2.91, 5.34] \\
    HGVRec $-$ +ASHA & 1.23 (2.91) [0.53, 1.92] & 1.22 (0.85) [0.03, 2.42] \\
    \bottomrule
  \end{tabular}
\end{table}

\subsection{Experimental Setup}
\label{sec:experimental-setup}
\paragraph{Datasets and protocol}
We use MicroLens-100K for its micro-video interactions and video content~\cite{ni2025microlens}, and M$^3$L-10M, which adds multimodal movie information, including trailers, to MovieLens-10M~\cite{spillo2026bingewatch}.
For M$^3$L-10M, every rating is treated as an interaction, movies without usable trailer videos are excluded, and the task is predicting the next movie in a rating-event sequence.
Users must have at least five interactions, and we retain their 13 most recent events; the candidate catalog is the union of items in these retained sequences.
This yields 100,000 users, 19,671 items, and 678,355 interactions for MicroLens-100K, and 69,878 users, 8,197 items, and 908,412 interactions for M$^3$L-10M.
The last event of each user is held out for test and the penultimate event for validation.
Validation selects the checkpoint with the highest HR@10; test uses it without refitting, ranking the full catalog after masking all items in the preceding history.

\paragraph{Representation variants}
All variants use the same content-based sequential recommender.
\emph{Title-only} projects the title embedding and applies a residual item MLP.
\emph{Scene Description} combines a projected title with projected scene-description embeddings, applies a residual scene MLP, and averages the scenes.
\emph{Scene Graph} projects title, entity, action, and context features, pools entities and actions separately within each scene, and averages the encoded scenes.
\emph{Scene Graph + ASHA} adds the action-segment module, and \emph{HGVRec} further adds REAF.
The three graph variants share identical extracted graphs, which are valid for 96.04\% and 95.92\% of the scenes in the two datasets; descriptions are a separate extraction.
Scene graphs and descriptions are extracted with Qwen3-VL and embedded with a frozen BGE text encoder.
We train with AdamW (weight decay 0.1, gradient-norm clipping 5.0) and seeds 42, 43, and 44.
% TODO(R05-config): R05 does not archive training.json or the extraction settings. Before
% submission, confirm and state the exact MLLM/BGE checkpoints (was the extractor the 2B
% Instruct model? "compact MLLM" in the Introduction depends on it), window length, frame
% count/resolution, token limits, prompt revisions, d, D, L, Transformer blocks/heads/dropout,
% learning rate, batch size, epoch limit, patience, precision and hardware.
% Do not copy the current config.yaml as R05 settings.

\paragraph{Metrics and statistics}
We report the mean and sample standard deviation of NDCG@10 and HR@10 over the three seeds, with relative changes over Title-only.
For NDCG@10, we average each user's metric over seeds and perform a paired bootstrap over users with 10,000 replicates; Bonferroni correction over all ten pairwise comparisons per dataset at family-wise level $0.05$ yields 99.5\% confidence intervals.
HR@10 is reported descriptively.

\subsection{Results}
\label{sec:overall-results}
\paragraph{Overall performance (RQ1)}
HGVRec achieves the highest mean NDCG@10 and HR@10 among the five variants on both datasets (Table~\ref{tab:overall-results}), with larger relative gains on MicroLens-100K.
Its NDCG@10 intervals against Title-only and Scene Description exclude zero on both datasets (Table~\ref{tab:contrasts}).
Scene Graph alone is below Scene Description on both datasets, indicating that producing graph-formatted outputs and pooling their components is insufficient by itself to reproduce the gains of HGVRec.

\paragraph{Effect of ASHA (RQ2)}
\label{sec:asha-results}
With identical graph inputs, adding ASHA to Scene Graph raises NDCG@10 by 20.14\% and 2.99\%; both corrected confidence intervals are positive, and mean HR@10 also increases.
The comparison adds the complete module (action-segment membership, size-aware segment encoding, scene updates, and segment readout) and its 1.05M trainable parameters.

\paragraph{Effect of REAF (RQ3)}
\label{sec:reaf-results}
Adding REAF to Scene Graph + ASHA increases NDCG@10 by a further 2.91\% and 0.85\% with 83K additional parameters, and both corrected intervals exclude zero.
This supports the incremental benefit of encoding actor, target, tool, and place assignments before pooling actions in the presence of ASHA.

% Working notes, not part of this submission (kept out of the text on purpose):
% matched scene-sequence and attention comparisons, stronger-MLLM extractors, a REAF-only
% variant and parameter-matched controls, grounding/explanation evaluation, and measured
% device latency/memory. Add claims about any of these only together with their results.
